% 原書pp.359--364、conifold transitionの項末まで。
\sectionlink*{用語集}
\newcommand{\gloss}[2]{%
  \par\medskip\begingroup
  \hangindent=1.5em\hangafter=1
  \noindent\textbf{#1}\quad #2\par
  \endgroup
}

\gloss{$(0,2)$コンパクト化（$(0,2)$ compactification）}
{$d=4,\ N=1$超対称性を持つ一般的なヘテロ弦の真空。世界面CFTは、右向きに伝播する$N=2$超共形不変性を持つ。}

\gloss{$(2,0)$理論（$(2,0)$ theory）}
{現在の用法では、張力のない弦と$d=6$の$(2,0)$超対称性を持つ、非自明な固定点理論の一群を指す。これらは、重なったM5ブレーンおよびIIAのNS5ブレーン上に、またA--D--E特異点におけるIIB理論に現れる。}

\gloss{$(2,2)$コンパクト化（$(2,2)$ compactification）}
{$d=4,\ N=1$のヘテロ弦真空のうち、Calabi--Yauコンパクト化を含む特別な部分集合に属するもの。世界面CFTは、右向きと左向きの両方に$N=2$超共形不変性を持つ。II型弦理論では、これらのCFTは$d=4,\ N=2$超対称性を持つ真空を与える。}

\gloss{A--D--E特異点（A--D--E singularity）}
{1つ以上の2次元球面が体積ゼロに収縮することによって生じる、実4次元の複素多様体の特異点。A--D--Eという用語は、収縮した球面の交点数を記述する、単純結線型Lie代数のDynkin図を指す。}

\gloss{Abel微分（Abelian differential）}
{Riemann面上で大域的に定義された正則$(1,0)$形式。}

\gloss{abstruse恒等式（abstruse identity）}
{Jacobiによるテータ関数の4次の恒等式の組の1つ。GSO射影を施した弦理論で、超対称性が要求するボソンとフェルミオンの縮退を導く。}

\gloss{アフィンLie代数（affine Lie algebra）}{カレント代数（current algebra）を参照。}

\gloss{アノマリー（anomaly）}
{量子効果によって古典的な対称性が破れること。\emph{重力アノマリー}は座標不変性のアノマリーである。\emph{大域的アノマリー}は、大きな対称変換、すなわち恒等変換に連続的につながっていない変換におけるアノマリーである。}

\gloss{アノマリー多項式（anomaly polynomial）}
{$d$次元における形式的な$(d+2)$形式であり、ゲージ・アノマリーと重力アノマリーの情報を表す。}

\gloss{漸近的局所Euclid空間（asymptotically locally Euclidean (ALE) space）}
{長距離では、離散群による同一視を施した平坦なEuclid空間に近づく空間。これは、オービフォールドの固定点、またはブローアップした固定点の近傍の幾何である。}

\gloss{Atiyah--Drinfeld--Hitchin--Manin（ADHM）構成}
{自己双対な場の強さを持つ、すべてのYang--Mills場の配位を構成する方法。}

\gloss{補助場（auxiliary field）}
{伝播しない場。すなわち、その場の方程式が微分方程式ではなく代数方程式である場。多くの超対称理論では、このような場を導入することによって変換を簡単にできる。}

\gloss{アクシオン（axion）}
{自発的に破れたPQ対称性に伴うGoldstoneボソン。\emph{模型に依存しないアクシオン}は、すべての摂動的な弦理論に現れ、重力子およびディラトンと密接に関係する。}

\gloss{$bc$ CFT}
{反可換な場からなる自由CFTで、微分について1階の作用を持つ。ウェイト$h_b=1-h_c$をパラメーターとする、このようなCFTの族がある。$h_b=2$のとき、このCFTは共形不変性に伴うFaddeev--Popovゴーストを記述する。}

\gloss{$\beta\gamma$ CFT}
{可換な場からなる自由CFTで、微分について1階の作用を持つ。ウェイト$h_\beta=1-h_\gamma$をパラメーターとする、このようなCFTの族がある。$h_\beta=\frac32$のとき、このCFTは超共形不変性に伴うFaddeev--Popovゴーストを記述する。}

\gloss{Batalin--Vilkovisky形式（Batalin--Vilkovisky formalism）}
{拘束のある、より一般的な理論を量子化するためのBRST形式の拡張。弦の場の理論で有用であった。}

\gloss{Becchi--Rouet--Stora--Tyutin（BRST）不変性}
{Faddeev--Popovの方法でゲージ固定された理論の冪零な対称性であり、もとのゲージ対称性に含まれる情報を表す。}

\gloss{Beltrami微分（Beltrami differential）}
{Riemann面の複素構造を、そのモジュライについて微分したもの。}

\gloss{Berezin積分（Berezin integration）}
{Grassmann変数の関数を複素数に移す線形演算であり、通常の積分の重要な性質の多くを持つ。}

\gloss{ベータ関数（beta function）}
{1. 量子場理論では、相互作用の有効な強さの長さスケールに関する微分。2. ガンマ関数の比で表され、Veneziano振幅に現れる特殊関数。}

\gloss{Betti数（Betti numbers）}
{de Rhamコホモロジーにおける非自明な$p$形式の数。$B_p$と表す。}

\gloss{ブラックホール・エントロピー（black hole entropy）}
{ブラックホールの地平面の面積に比例する量$S$で、$S=2\pi A/\kappa^2$。これは熱力学的エントロピーの性質を持つ。古典的一般相対論では減少せず、Hawking放射を含めても、ブラックホール・エントロピーと通常のエントロピーの和は減少しない。このエントロピーを統計力学的に導くことは大きな目標であり、最近の研究で部分的に実現された。}

\gloss{ブラックホール蒸発（black hole evaporation）}
{地平面の近傍での対生成によって、ブラックホールが熱放射、すなわちHawking放射を放出すること。}

\gloss{ブラックホール情報パラドックス（black hole information paradox）}
{量子力学と一般相対論の間の矛盾。ブラックホールに落ち込んだ情報は失われ、ブラックホールが蒸発しても再び現れない。これは通常の量子力学的な時間発展と両立しない。どうやら、量子力学を大きく修正するか、局所性についての通常の理解が大きく破綻するかの、いずれかが必要となる。}

\gloss{ブラック$p$ブレーン（black $p$-brane）}
{事象の地平面を持つ$p$次元の広がった物体。$p$方向には並進不変で、残りの方向にはブラックホールの幾何を持つ空間である。}

\gloss{ブローアップ（blow up）}
{特異な多様体を滑らかな多様体へ変形すること。}

\gloss{Bogomolnyi--Prasad--Sommerfield（BPS）状態}
{超対称代数全体の非自明な部分代数のもとで不変な状態。このような状態は常に保存電荷を持ち、超対称代数はその電荷を用いて状態の質量を厳密に決定する。BPS状態は、非BPS状態よりも小さい超対称性の表現、いわゆる\emph{短い表現}に属する。大きさの異なる短い表現がある場合には、可能な最小の表現である\emph{超短い表現}も区別する。一般に、その次元は非BPS表現の次元の平方根である。}

\gloss{Borel総和（Borel summation）}
{発散級数の和を定義する方法。場の理論や弦理論の非摂動効果を研究する手段として使われてきた。ただし、ほとんどの非摂動効果は摂動級数を用いて有効に研究できるものではないことを理解しておくべきである。}

\gloss{Born--Infeld作用（Born--Infeld action）}
{通常のゲージ場の作用の一般化で、ゲージ場の強さについて非多項式的である。もともとは電磁気学を短距離で修正する可能性として提案された。Dブレーン上のゲージ場の低エネルギー有効作用として現れる。}

\gloss{ボソン化（bosonization）}
{フェルミオン場の理論とボソン場の理論との厳密な同値性。2次元で可能である。ボソンはフェルミオンと反フェルミオンの対であり、フェルミオンはボソンのコヒーレント状態である。}

\gloss{$c$写像（$c$-map）}
{Calabi--Yau 3次元複素多様体にコンパクト化した一方のII型弦理論のハイパー多重項モジュライ空間を、同じ多様体上の他方のII型理論のベクトル多重項モジュライ空間から構成する方法。}

\gloss{$c$定理（$c$-theorem）}
{2次元のユニタリーなCFTにおいて、長さスケールが増すと単調非増加であり、固定点では停留して中心電荷に等しくなる正の量$c$が存在すること。これは繰り込み群の流れの大域的な形に強い制約を与える。高次元には単純な類似物は存在しないように思われる。Zamolodchikovの$c$定理とも呼ばれる。}

\gloss{$CPT$対称性（$CPT$ symmetry）}
{パリティ反転、時間反転、荷電共役を組み合わせた操作。すべてのLorentz不変な局所量子場理論の対称性である。}

\gloss{Calabi--Yau多様体（Calabi--Yau manifold）}
{第1 Chern類が消えるKähler多様体。Calabi--Yau $n$次元複素多様体は、$2n$個の実座標、すなわち$n$個の複素座標を持つ。Yauの定理によって、$SU(n)$ホロノミーを持つRicci平坦な計量の存在が保証される。}

\gloss{伝播関数の相殺による議論（canceled propagator argument）}
{広い条件のもとで、Riemann面のモジュライ空間上の表面項が消え、したがって弦の振幅から非物理的な状態が分離することを導く一般原理。このような振幅は、境界で被積分関数が急速に減衰し、表面項が恒等的にゼロになる領域からの解析接続によって定義される。したがって、その解析接続も恒等的にゼロである。}

\gloss{Casimirエネルギー（Casimir energy）}
{場に課した境界条件に起因する、量子場理論の基底状態エネルギーの変化。}

\gloss{重心モード（center-of-mass mode）}
{量子場の空間的なFourier展開の第0成分。}

\gloss{中心電荷（central charge）}
{Lie代数の右辺に現れ、その代数のすべての演算子と可換な演算子。定数である場合もある。代表的な例には、Virasoro代数の定数項と、多くの超対称代数の右辺に現れる電荷がある。}

\gloss{Chan--Paton自由度（Chan--Paton degrees of freedom）}
{開弦の端点に局在する自由度。現在では、弦が終端するDブレーンを指定するものと解釈される。}

\gloss{Chan--Paton因子（Chan--Paton factor）}
{Chan--Paton自由度の状態に対応する頂点演算子の因子。}

\gloss{Chern--Simons項（Chern--Simons term）}
{場の強さに加えて$p$形式ポテンシャルを含む作用の項。このような項は、Bianchi恒等式、および／または$p$形式のゲージ変換の修正によって、ゲージ不変である。通常、これらの項は位相やアノマリーと密接に関係する。}

\gloss{カイラル（chiral）}
{1. パリティについて非対称に作用すること。カイラル多重項、カイラル対称性、カイラル理論、カイラリティ、拡張カイラル代数を参照。2. 超対称代数の一部のもとで不変であること。カイラル場、カイラル多重項、カイラル一次場、カイラル環を参照。}

\gloss{カイラル場（chiral field）}
{超対称性において、代数の一部のもとで不変な局所演算子。BPS状態に対応する演算子版である。}

\gloss{カイラル多重項（chiral multiplet）}
{2つの実スカラーを持つ、$d=4,\ N=1$超対称性の多重項。クォークとレプトンはこのような多重項に含まれる。この多重項は、カイラルという語の両方の意味に関係する。すなわち、カイラルな結合を持つフェルミオンを含み、対応する超場の積分が超対称代数の半分のもとで不変である。}

\gloss{カイラル一次場（chiral primary）}
{$N=2$ SCFTにおいて、剛体超対称性$G^\pm_{-1/2}$の一方によっても消される一次場。}

\gloss{カイラル環（chiral ring）}
{カイラル場の閉じたOPE代数。}

\gloss{カイラル対称性（chiral symmetry）}
{スピノル場への作用がパリティ対称でない対称性。}

\gloss{カイラル理論（chiral theory）}
{ゲージ結合がパリティ対称でないゲージ理論。}

\gloss{カイラリティ（chirality）}
{$d=2k$次元において、$\Gamma^\mu$と反可換な演算子$\Gamma^d$の固有値。この固有値によって、パリティで互いに移り合う2つのWeyl表現が区別される。}

\gloss{Christoffel接続（Christoffel connection）}
{一般相対論において、計量から構成される接続。}

\gloss{臨界挙動（critical behavior）}
{質量のない場を持つ赤外固定点での量子場理論の振舞いと、その振舞いへの近づき方。}

\gloss{閉弦（closed string）}{円の位相を持つ弦。}

\gloss{コサイクル（cocycle）}
{頂点演算子において、生成・消滅演算子について正規順序化された指数関数に掛かる、演算子値の位相因子。演算子に正しい交換または反交換の性質を与えるため、これが必要となる場合がある。}

\gloss{係数関数（coefficient functions）}
{演算子積の展開に現れる局所演算子に掛かる、位置に依存する係数。}

\gloss{コホモロジー（cohomology）}
{冪零演算子$Q$、すなわち$Q^2=0$を満たす演算子を持つ任意のベクトル空間において、$Q$の核を$Q$の像で割ったもの。つまり、閉じた状態（$Q$によって消される状態）の空間で、完全な状態（$Q\psi$の形の状態）をゼロと同値と定めたものである。\emph{de Rhamコホモロジー}は、微分形式に作用する外微分$d$のコホモロジーである。複素多様体上の\emph{Dolbeaultコホモロジー}は、$(p,q)$形式上の$\partial$と$\bar\partial$、すなわち$d$の$(1,0)$部分と$(0,1)$部分のコホモロジーである。\emph{ホモロジー}は境界演算子のコホモロジーである。\emph{BRSTコホモロジー}はBRST演算子のコホモロジーであり、ゲージ不変な理論の物理的な状態空間を定義する。}

\gloss{Coleman--Weinberg公式（Coleman--Weinberg formula）}
{自由な量子場の真空エネルギー密度を、そのモードの零点エネルギーの繰り込まれた和から与える式。}

\gloss{収縮するサイクル（collapsing cycle）}
{ある極限で体積が消えるサイクル。通常は特異な多様体を生じる。}

\gloss{集団座標（collective coordinate）}
{ソリトンやその他の広がった物体を量子化するとき、その位置や配位に対応する自由度。}

\gloss{コンパクトなCFT（compact CFT）}
{任意に与えられた値よりエネルギーの小さい状態の数が有限であるCFT。コンパクト空間上の微分演算子のスペクトルとの類似によって定義される。}

\gloss{コンパクト化スケール（compactification scale）}
{波動関数がコンパクトな次元に非自明に依存する状態の、特徴的な質量スケール。}

\gloss{コンパクト化する（compactify）}
{空間方向の次元の一部がコンパクトである時空で、場の理論や弦理論を考えること。}

\gloss{複素多様体（complex manifold）}
{複素座標系が与えられ、その座標の正則な再パラメーター化を同一視した多様体。}

\gloss{複素構造（complex structure）}
{複素座標の同値類。与えられた微分可能多様体には、互いに同値でない多くの複素構造が存在しうる。}

\gloss{複素構造モジュライ（complex structure moduli）}
{多様体上の同値でない複素構造をパラメーター化するモジュライ。Calabi--Yau 3次元複素多様体へのコンパクト化では、$(2,1)$形式に対応する。}

\gloss{共形ブロック（conformal block）}
{CFTにおいて、状態についての和に対する、1つの共形族からの寄与。}

\gloss{共形ブートストラップ（conformal bootstrap）}
{対称性と整合性の条件だけを用いてすべてのCFTを構成するという、部分的に成功した研究計画。}

\gloss{共形族（conformal family）}
{最高ウェイト状態にVirasoroの上昇生成子を、互いに同値でないあらゆる仕方で作用させて得られる状態の集合。または、対応する局所演算子の集合。\emph{縮退した共形族}は、その族のすべての状態に直交する零状態を含む。}

\gloss{共形場理論（conformal field theory, CFT）}
{共形不変な量子場理論。}

\gloss{共形ゲージ（conformal gauge）}
{計量が単位計量に比例するような、2次元での座標の選び方。}

\gloss{共形Killingベクトル（conformal Killing vector）}
{計量を不変にする、大域的に定義された無限小の$\mathrm{diff}\times\mathrm{Weyl}$変換。}

\gloss{共形変換（conformal transformation）}
{Euclid空間またはMinkowski空間をそれ自身に写し、位置に依存する再スケールを除いて平坦な計量を不変にする写像。同値な言い方をすれば、平坦な計量を不変に保つ$\mathrm{diff}\times\mathrm{Weyl}$の部分群。$d\geq3$次元では$\frac12(d+1)(d+2)$個のパラメーターを持つ。2次元では、すべての正則写像の集合である。有限変換を考えるには、球面のMöbius変換の場合と同様に、無限遠点を含める必要がある。

この用法は弦理論と量子場理論では標準となっている。しかし一般相対論では、共形変換は、位置に依存する計量の任意の再スケールと定義される。これは現在、弦理論ではWeyl変換と呼ばれている。}

\gloss{コニフォールド（conifold）}
{3サイクルの収縮に対応する、特異な複素構造を持つCalabi--Yau多様体。この空間上の弦理論は特異である。現在では、この特異性は、サイクルに巻き付いた質量のない3ブレーンの量子効果に起因すると理解されている。}

\gloss{コニフォールド転移（conifold transition）}
{質量のない3ブレーン場の凝縮に起因する位相の変化。適切な条件のもとでは、複数のサイクルが収縮したコニフォールド上の、質量のない3ブレーン場のポテンシャルは、これらの場について平坦な方向を持つ。これは、収縮した3サイクルの代わりに2サイクルをブローアップするという、位相の変化に対応する。}

